% ======================================================================
\section{Electronic models}
\label{sec:models}
% ======================================================================

\subsection{Doped monolayer graphene}
\label{sec:slgmodel}

Near one Dirac point, $h(\bp)=\hbar v_F(\tau p_x\sigma_x+p_y\sigma_y)$ with $\bp$ measured from the Dirac
point, $\hbar v_F=0.658$ eV\,nm ($v_F=10^6$ m/s) and $g=g_sg_v=4$. The bands and spinors are
\begin{equation}
\xi^s_{\bp}=s\hbar v_F|\bp|-E_F,\qquad u^s_{\bp}=\tfrac{1}{\sqrt2}(1,\,s\,e^{\ii\phi_{\bp}}),\qquad s=\pm1,
\end{equation}
where $\tau=\pm1$ labels the valley, $\phi_{\bp}$ is the polar angle of $\bp$, $k_F=\sqrt{4\pi n/g}$ the Fermi momentum at
electron density $n$ and $E_F=\hbar v_Fk_F$ the Fermi energy ($n=10^{12}$ cm$^{-2}$ $\Rightarrow$ $k_F=0.1772$ nm$^{-1}$,
$E_F=116.6$ meV). The overlap that dresses every Coulomb vertex is the chirality factor
\begin{equation}
F^{s's}(\bp,\bp')=|\braket{u^{s'}_{\bp}}{u^{s}_{\bp'}}|^2=\tfrac12\bigl[1+ss'\cos\angle(\bp,\bp')\bigr],
\qquad \textstyle\sum_sF^{s's}=1,
\label{eq:chirality}
\end{equation}
which forbids intraband backscattering and interband forward scattering
(\code{slg\_self\_energy.make\_dirac\_form\_factor})~\cite{CastroNeto2009,Kotov2012}. The sample has only one sheet,
so the layer-projected form factor of \cref{sec:formfactors} is $\rho^\ell=\delta_{\ell S}\braket{u}{u'}$, where $S$ labels
the sample sheet: a sample electron has no intermediate state in the tip.

\subsection{Twisted bilayer graphene: continuum model}
\label{sec:tbgmodel}

\paragraph{Hamiltonian.} One valley ($\tau=+$) of the Bistritzer--MacDonald model in the form of
Koshino \emph{et al.}~\cite{BM2011,Koshino2018} as implemented in \texttt{pybandstructure}
(\code{koshino}). In the basis of plane waves $\bk+\bG$ ($\bG$ in the moir\'e reciprocal lattice,
$|\bG|\le G_{\max}|\bG_M|$ with $G_{\max}=3$, i.e.\ 37 plane waves, so that $H(\bk)$ is a $148\times148$ matrix) and of the
four orbitals
$(\ell,\sigma)\in\{t,b\}\times\{A,B\}$,
\begin{gather}
H_{\bG\bG'}(\bk)=\delta_{\bG\bG'}
\begin{pmatrix}h_t(\bk+\bG)&0\\0&h_b(\bk+\bG)\end{pmatrix}
+\sum_{j=1}^3\Bigl[\delta_{\bG-\bG',\bm g_j}\begin{pmatrix}0&T_j\\0&0\end{pmatrix}+{\rm h.c.}\Bigr],
\nonumber\\
h_\ell(\bk)=\hbar v\,\bigl[R_\ell(\bk-\bK_\ell)\bigr]\cdot(\tau\sigma_x,-\sigma_y),
\end{gather}
where $R_{t,b}=O(\pm\theta_{\rm TBG}/2)$ rotates the two layers in opposite directions, and
$\bK_{1,2}=\frac{|\bG_M|}{\sqrt3}(-\frac{\sqrt3}{2},\pm\frac12)$ are the moir\'e $\kappa$ points at which the top ($1$) and bottom
($2$) Dirac cones sit. The moir\'e reciprocal lattice is spanned by $\bm b_1,\bm b_2$ with
$|\bm b_{1,2}|=|\bG_M|=8\pi\sin(\theta_{\rm TBG}/2)/(\sqrt3a)$, $a=0.246$ nm the graphene lattice constant. The interlayer
tunneling connects plane waves that differ by $\bm g_j\in\{0,\bm b_1,\bm b_2\}$ with the $2\times2$ sublattice matrices of
Koshino \emph{et al.}~\cite{Koshino2018}
\begin{equation}
T_1=\begin{pmatrix}u_{AA}&u_{AB}\\u_{AB}&u_{AA}\end{pmatrix},\quad
T_2=\begin{pmatrix}u_{AA}&u_{AB}\varpi^{-\tau}\\u_{AB}\varpi^{\tau}&u_{AA}\end{pmatrix},\quad
T_3=\begin{pmatrix}u_{AA}&u_{AB}\varpi^{\tau}\\u_{AB}\varpi^{-\tau}&u_{AA}\end{pmatrix},\qquad\varpi=e^{2\pi\ii/3},
\end{equation}
where $u_{AA}$ and $u_{AB}$ are the tunneling amplitudes between equal and opposite sublattices (the assignment of $T_j$ to
$\bm g_j$ follows \code{pybandstructure.koshino}). Parameters (Ref.~\cite{Wei2025}):
$\theta_{\rm TBG}=1.20^\circ$, $u_{AB}=117$ meV, $u_{AA}=0.7u_{AB}$ (the reduced value mimics lattice relaxation) and
$\hbar v=0.657$ eV\,nm. The two central (flat) bands are then 24--25 meV wide and separated from the remote bands by
$\approx44$ meV. The code also reports the Bistritzer--MacDonald coupling $\alpha=u_{AB}/(2\hbar v_0k_D\sin\frac{\theta_{\rm TBG}}2)$
with the reference value $\hbar v_0k_D=9.905$ eV of Ref.~\cite{BM2011} (\code{TBGParams.alpha}, $\alpha=0.56$), and the
corresponding first-order estimate of the renormalised Dirac velocity~\cite{BM2011}
\begin{equation}
v_{\rm BM}=v\,\frac{1-3\alpha^2}{1+6\alpha^2}
\label{eq:vstar}
\end{equation}
($\hbar v_{\rm BM}=0.010$ eV\,nm, \code{TBGParams.v\_renorm}). Both numbers are used only to set the default broadening $\eta$
(\cref{sec:numerics}). With the actual $\hbar v$ the coupling would be $u_{AB}/(\hbar v|\bK_1-\bK_2|)=0.50$, so $v_{\rm BM}$ is an order
of magnitude estimate of the flat-band velocity, not a fitted value. Only the $2n_b=8$ bands closest to neutrality are computed
(\code{TBGParams.n\_bands}$=4$). The electrostatic environment differs from Ref.~\cite{Wei2025}, which uses an
isotropic background $\varepsilon_r=10$: here the TBG sits in anisotropic hBN,
$(\varepsilon_\parallel,\varepsilon_z)=(6.41,3)$, with vacuum between its layers (\code{TBGSLGConfig.eps\_par},
\code{eps\_z}, \code{eps\_interlayer}), so plasmon energies are not expected to coincide with that reference.

\paragraph{Model frame vs lab frame.} The model momentum $\bk$ is measured from the moir\'e $\gamma$ point.
The graphene (lab) momentum of the plane wave $\bk+\bG$ in either layer is
\begin{equation}
\bp=\bk+\bG+\bm\Gamma^{\rm lab},\qquad \bK^{\rm lab}_{\rm top}=\bm\Gamma^{\rm lab}+\bK_1,\quad
\bK^{\rm lab}_{\rm bot}=\bm\Gamma^{\rm lab}+\bK_2,
\label{eq:labframe}
\end{equation}
where $\bm\Gamma^{\rm lab}$ (\code{gamma\_point\_tblg}) lies on the $x$ axis and $|\bK^{\rm lab}_{\rm top,bot}|=4\pi/3a$
at angles $\pm\theta_{\rm TBG}/2$ (\code{tbg\_model.model\_frame}). This is checked in
\code{test\_model\_frame\_dirac\_points}, where the decoupled model has a zero-energy top-layer state at $\bK_1$.
This relation is what the QTM kinematics needs (\cref{sec:qtmtbg}).

\paragraph{Bloch states.} $H(\bk)u_{n\bk}=E_{n\bk}u_{n\bk}$, where the eigenvector has components
$u_{n\bk}(\bG,\ell,\sigma)$, stored as \code{u[4 i\_G + 2l + sigma]} with $i_{\bG}$ the index of the plane wave $\bG$. The states are stored once
(\code{band\_structure.wavefunctions}) and the \emph{same} gauge is used for the form factors, the self-energy
and the tunneling amplitudes. The chemical potential is fixed by the filling at temperature $T=5$ K
(charge neutrality by default), and energies are measured from it, $\xi_{n\bk}=E_{n\bk}-\mu$.

\paragraph{Layer-resolved form factors.} The Coulomb vertices of the screening and of the self-energy are the matrix
elements of the density of each layer between Bloch states, $\rho^\ell_{nm}(\bk,\bq)=\bra{u_{n\bk}}P_\ell\,\hat\tau_{\bG_0}\ket{u_{m,\bk+\bq}}$
(\code{densityL[l, n, m, k, q, G]}), where $P_\ell$ projects on layer $\ell$ and $\hat\tau_{\bG_0}$ compensates the folding of
$\bk+\bq$ into the first zone. Their derivation from the interaction Hamiltonian, and the reason why the layer index is
needed, are given in \cref{sec:hint,sec:formfactors}.

\paragraph{Doping and Hartree.} For a doped sample the first-shell Hartree potential
$V_H=\sum_{|\bG|=|\bG_M|}\frac{2\pi e^2}{\varepsilon_H|\bG|}\delta n(\bG)e^{\ii\bG\cdot\br}$, with $\varepsilon_H$ the background
dielectric constant and $\delta n(\bG)$ the Fourier components of the electron density, can be solved
self-consistently (\code{tbg\_model.apply\_hartree}, which needs the development version of \texttt{pybandstructure}).
The densities $\delta n(\bG)$ are measured from charge neutrality (regularised occupations), and
$G_0$ is then built on the Hartree bands. At neutrality the Hartree potential of the BM model vanishes.

\subsection{Moir\'e Brillouin-zone meshes}

$\bk$ and $\bq$ run over the same mesh of $N_s\times N_s$ points of the moir\'e Brillouin zone, mapped into its
Wigner--Seitz cell (\code{n\_sites}$\,=N_s$, \cref{fig:quadrature}a). Every mesh point represents the same fraction
$1/N_s^2$ of the zone, whose area is $A_{\rm BZ}=\frac{\sqrt3}{2}|\bG_M|^2$. Momentum integrals become sums with the
weights $w_{\bq}=A_{\rm BZ}/(2\pi N_s)^2$, derived in \cref{eq:qweights}. These weights do not contain the spin--valley
degeneracy, which is added explicitly where it belongs (in $\chi^0$, \cref{sec:degeneracy}). The band structure is
computed for valley $K$, and valley $K'$ follows by time reversal (\cref{sec:valleys}).
